\documentclass[11pt]{article}

\usepackage[margin=1in]{geometry}
\usepackage[T1]{fontenc}
\usepackage[utf8]{inputenc}
\usepackage{lmodern}
\usepackage{xcolor}
\usepackage{hyperref}
\usepackage{enumitem}
\usepackage{array}
\usepackage{tabularx}
\usepackage{fancyhdr}
\usepackage{titlesec}

\hypersetup{
  colorlinks=true,
  linkcolor=black,
  urlcolor=blue,
  pdftitle={Beginner Aquaponics Handout},
  pdfauthor={OpenAI},
  pdfsubject={Introductory aquaponics handout},
  pdfkeywords={aquaponics, media bed, beginner, fish tank, grow bed}
}

\setlength{\parindent}{0pt}
\setlength{\parskip}{0.65em}
\setlist[itemize]{leftmargin=1.5em}
\setlist[enumerate]{leftmargin=1.75em}

\pagestyle{fancy}
\fancyhf{}
\lhead{Beginner Aquaponics Handout}
\rhead{\thepage}
\renewcommand{\headrulewidth}{0.4pt}

\titleformat{\section}{\large\bfseries\color{black}}{\thesection}{0.6em}{}
\titleformat{\subsection}{\normalsize\bfseries\color{black}}{\thesubsection}{0.6em}{}

\begin{document}

\begin{center}
  {\LARGE \textbf{Beginner Aquaponics Handout}}\\[0.5em]
  {\large A practical first-build guide for a small media-bed system}\\[0.75em]
  \rule{0.9\linewidth}{0.6pt}
\end{center}

\section*{Purpose}
This handout explains how to plan, build, and start a basic aquaponics system for learning or small-scale home production. It is written for beginners and emphasizes a conservative startup approach: keep the design simple, stock lightly, monitor water quality, and let the biological filter mature before expanding the system.

\section*{What Aquaponics Is}
Aquaponics combines fish culture and plant production in a recirculating water loop. Fish produce waste, nitrifying bacteria convert toxic ammonia into nitrite and then nitrate, plants use those nutrients, and cleaned water returns to the fish tank [1][2][3].

A beginner system is best understood as three living parts that must remain in balance:
\begin{itemize}
  \item \textbf{Fish}, which generate nutrients and require stable water quality.
  \item \textbf{Plants}, which remove nutrients and return cleaner water to the system.
  \item \textbf{Beneficial bacteria}, which perform nitrification and make the whole loop work.
\end{itemize}

\section*{Recommended First System}
For a first build, a \textbf{media-bed aquaponics system} is usually the most forgiving design. In a media-bed system, the grow media supports the plants and also provides a large amount of surface area for nitrifying bacteria [1][2]. A small DIY system commonly uses a food-grade IBC tote or drum as the fish tank, a grow bed positioned above or beside the tank, a water pump, a drain, and an air pump with an air stone [2].

\textbf{Why start with a media bed?}
\begin{itemize}
  \item It is structurally simple.
  \item It supports both plants and biofiltration.
  \item It is usually easier for beginners than systems that require separate solids filtration and more precise hydraulic tuning.
\end{itemize}

\section*{Core Components}
A basic beginner system should include the following:
\begin{itemize}
  \item \textbf{Fish tank}: a clean, food-grade container with known prior contents.
  \item \textbf{Grow bed}: placed so water can drain back to the tank.
  \item \textbf{Grow media}: inert media such as expanded clay, expanded shale, or suitable washed gravel [1][2].
  \item \textbf{Water pump}: sized to circulate water reliably.
  \item \textbf{Drain plumbing}: a standpipe, timer-based drain arrangement, or bell siphon.
  \item \textbf{Aeration}: an air pump and air stone to maintain dissolved oxygen.
  \item \textbf{Water test kit}: for ammonia, nitrite, nitrate, and pH at minimum [3].
\end{itemize}

\textbf{Important note:} a bell siphon is common in flood-and-drain media-bed systems, but it is not mandatory. A timed flood-and-drain approach can also work [1].

\section*{Build Procedure}
\subsection*{1. Choose safe containers}
Use a tank and grow bed made from materials suitable for food contact or known-safe horticultural use. Repurposed IBC totes and drums are common, but only when prior contents are known and the container is thoroughly cleaned [2].

\subsection*{2. Set the tank and grow bed}
Place the grow bed where gravity can return water to the fish tank. Make sure the structure can support the full weight of wet media, water, and plants.

\subsection*{3. Wash and add media}
Rinse the media thoroughly before filling the bed. This removes fine particles that would otherwise cloud the water and clog plumbing [1]. Media depth around 12--14 inches is commonly recommended for media beds [1].

\subsection*{4. Install plumbing}
Install the pump line from the fish tank to the grow bed, then install the drain back to the tank. In media-bed systems, keep the flood level slightly below the top of the media to reduce algae growth and limit stem or crown rot [4].

\subsection*{5. Add aeration}
Run an air pump with one or more air stones in the fish tank. Dissolved oxygen is important for fish health and also supports nitrifying bacteria [3].

\subsection*{6. Fill and test the system}
Fill the system with dechlorinated water, run the pump, and inspect the entire loop for leaks, unstable fittings, poor drainage, and low-flow areas.

\section*{Cycling the System}
A newly assembled system is not immediately ready for normal fish stocking. The biological filter must first develop bacteria that can process ammonia safely. This startup period is called \textbf{cycling} [4].

\textbf{What to expect during cycling}
\begin{enumerate}
  \item Ammonia appears first.
  \item Nitrite rises as bacteria begin converting ammonia.
  \item Nitrate appears once the second bacterial step is established.
  \item Ammonia and nitrite should both fall to low levels before heavier stocking is considered.
\end{enumerate}

Extension guidance for cycling with fish recommends advancing cautiously and using water testing rather than the calendar alone. A practical milestone is the presence of nitrate with ammonia and nitrite near or below about 0.5 ppm before increasing fish load [4].

\section*{Good Starter Fish and Plants}
\subsection*{Fish}
Common beginner choices include tilapia, catfish, and ornamental fish such as goldfish, depending on local law, water temperature, and whether the system is intended for food production or demonstration [1].

\subsection*{Plants}
Start with lower-demand crops such as leafy greens and herbs. Lettuce, basil, chives, spinach, and similar plants are generally better first crops than heavy fruiting plants such as tomatoes, cucumbers, or peppers [3].

\section*{Conservative Beginner Stocking}
Avoid aggressive stocking formulas during startup. A conservative reference point for a beginner media-bed system is a fish harvest weight of about \textbf{1 pound of fish per 8--10 gallons of water}, with light feeding and regular water testing used to decide when the system can support more biomass [4].

This is safer for a first build than relying on a simplified fish-to-media rule alone.

\section*{Routine Checks}
\textbf{Check daily:}
\begin{itemize}
  \item water circulation
  \item aeration
  \item fish behavior
  \item leaks or overflows
  \item clogged drains or inflow lines
  \item unusual cloudiness, solids buildup, or odor
\end{itemize}

\textbf{Check regularly with a test kit:}
\begin{itemize}
  \item ammonia
  \item nitrite
  \item nitrate
  \item pH
  \item temperature
  \item dissolved oxygen if available [3]
\end{itemize}

\section*{Common Beginner Mistakes}
\begin{itemize}
  \item adding too many fish too quickly
  \item overfeeding
  \item skipping aeration
  \item failing to test water during startup
  \item treating aquaponics like ordinary hydroponics instead of a living biofilter
  \item underestimating the need for daily inspection and maintenance [4][5]
\end{itemize}

It is also normal to need nutrient supplementation over time. Small aquaponic systems commonly require additions such as iron, potassium, and calcium because fish waste alone may not meet all plant nutrient needs [1].

\section*{Success Criteria for a First System}
Your first system is on the right track when:
\begin{itemize}
  \item water circulates reliably without leaks or flooding,
  \item fish are active and not gasping at the surface,
  \item ammonia and nitrite remain low,
  \item nitrate is present,
  \item plants show healthy new growth,
  \item the system can run day to day without constant manual correction.
\end{itemize}

\section*{Starter Checklist}
\begin{tabularx}{\linewidth}{>{\bfseries}l X}
Tank & Food-grade fish tank or known-safe container \\
Grow bed & Stable bed positioned to drain back to the tank \\
Media & Washed inert media \\
Pump & Water pump sized for reliable circulation \\
Drainage & Standpipe, timer drain, or bell siphon \\
Aeration & Air pump and air stone \\
Water & Dechlorinated water \\
Testing & Test kit for ammonia, nitrite, nitrate, and pH \\
Fish & Light initial stocking only \\
Plants & Easy starter crops such as lettuce or basil \\
\end{tabularx}

\section*{Selected Sources}
\begin{enumerate}[label={[\arabic*]}]
  \item Oklahoma State University Extension. \emph{Principles of Small-Scale Aquaponics}. \url{https://extension.okstate.edu/fact-sheets/principles-of-small-scale-aquaponics}
  \item University of Maryland Extension. \emph{How to Build a Low-Cost Small-Scale Aquaponic System (EM-2023-0698)}. \url{https://extension.umd.edu/resource/how-build-low-cost-small-scale-aquaponic-system-em-2023-0698}
  \item New Mexico State University. \emph{An Introduction to Aquaponics}. \url{https://pubs.nmsu.edu/_h/H170/index.html}
  \item Oklahoma State University Extension. \emph{Nitrification and Maintenance in Media-Bed Aquaponics}. \url{https://extension.okstate.edu/fact-sheets/nitrification-and-maintenance-in-media-bed-aquaponics}
  \item Oklahoma State University Extension. \emph{Key Facts for Prospective Aquaponics Producers}. \url{https://extension.okstate.edu/fact-sheets/key-facts-for-prospective-aquaponics-producers}
\end{enumerate}

\vfill
\hrule
\smallskip
{\footnotesize This handout is designed as a practical beginner reference. Local regulations, species requirements, and climate conditions should always be checked before stocking fish or scaling a system.}

\end{document}
